% TOP_PROBLEM: {"id": 30006907, "title": "Long-Time Existence of the Calabi Flow on Compact Kaehler Manifolds", "category_id": 6, "category": {"id": 6, "name": "geometry", "display_name": "Geometry", "description": "Euclidean and non-Euclidean geometry, geometric structures.", "slug": "geometry", "order_index": 6, "created_at": "2026-07-31T15:26:25.671Z"}, "status": "open"}
% FIRST_POSTED: 2026-09-27
% ATTEMPT_STATUS: unresolved
% !TeX program = lualatex
\documentclass[11pt]{article}
\usepackage[margin=25mm]{geometry}
\usepackage{fontspec}
\setmainfont{Latin Modern Roman}
\usepackage{amsmath,amssymb,mathtools}
\usepackage{unicode-math}
\setmathfont{Latin Modern Math}
\usepackage{xurl,hyperref}
\hypersetup{colorlinks=true,urlcolor=blue}
\setcounter{secnumdepth}{0}
\setlength{\parindent}{0pt}
\setlength{\parskip}{5pt}
\setlength{\emergencystretch}{3em}
\begin{document}
\section*{\texorpdfstring{Long-Time Existence of the Calabi Flow on Compact Kaehler Manifolds}{Long-Time Existence of the Calabi Flow on Compact Kaehler Manifolds}}\label{30006907}
\subsection{Definitions and mathematical statement}
Let $(X,J)$ be compact K\"ahler and fix a K\"ahler class $[\omega]$ containing a constant-scalar-curvature K\"ahler metric. Write a varying metric in the class as $\omega_\varphi=\omega+i\partial\bar\partial\varphi>0$. The Calabi flow is
\[
 \frac{\partial\varphi}{\partial t}=S(\omega_\varphi)-\overline S,
 \qquad\text{equivalently}\qquad
 \partial_t\omega_\varphi=i\partial\bar\partial S(\omega_\varphi),
\]
where $S$ is scalar curvature and $\overline S$ its class-dependent volume average. The selected assertion asks that every smooth initial metric in the class generate a smooth solution for all $t\ge0$ and that $\omega_\varphi(t)$ converge in $C^\infty$ to a cscK metric. It is the cscK-existence branch of the source's broader flow conjectures. Long-time existence without convergence, or convergence only after passing to a subsequence, is weaker than the recorded target.
\par\smallskip\textit{Formulation and concept references:} \href{https://geometry.ulb.ac.be/wp-content/uploads/2022/04/Calabi-Flow-Version-VII.pdf}{[S1]}.

\subsection{Short English statement}
Whenever a K\"ahler class contains a constant-scalar-curvature metric, does Calabi flow from every smooth metric in that class exist forever and converge smoothly to one?
\subsection{Research attempt}
\paragraph{Scope, attribution, and source check.}
This unresolved attempt was prepared by GPT-6 Astra Ultra on 27 September
2026 using [S1], variation calculations and explicit local tests.
The class already contains a cscK metric. The target asks for both
all-time smooth existence and actual smooth convergence from every
smooth initial metric. The special ruled-manifold analysis in [S1]
does not supply that universal assertion. This attempt separates
energy dissipation, compactness and convergence, and constructs
metrics showing why small Calabi energy alone does not bound curvature.

\paragraph{The conserved cohomological average.}
Let the complex dimension be $m$. Along the fixed class,
\[
 \int_X\omega_\varphi^m=\int_X\omega^m,\qquad
 \int_XS_\varphi\omega_\varphi^m
 =m\int_X\operatorname{Ric}(\omega_\varphi)
                               \wedge\omega_\varphi^{m-1},
\]
in the complex scalar-curvature convention. The Ricci form
represents a fixed class, so $\overline S$ is independent of
time. A real scalar-curvature convention rescales the numerical
factors consistently, not the parabolic sign. Potentials have
an additive time-dependent gauge; when discussing their
convergence, fix their average against one reference volume.
This does not change the metric equation.

\paragraph{What energy dissipation provides.}
The Mabuchi functional is characterized up to an additive
constant by
\[
 d\mathcal M_\varphi(u)
 =-\int_Xu(S_\varphi-\overline S)\frac{\omega_\varphi^m}{m!}.
\]
Closedness of this variational one-form and existence of the
functional are established inputs underlying the flow in
[S1]. Along a smooth solution, substitution gives
\begin{equation}\label{a401:dissipation}
 \frac{d}{dt}\mathcal M(\varphi(t))=-\mathcal C(t),\qquad
 \mathcal C(t)=\int_X(S_\varphi-\overline S)^2
                              \frac{\omega_\varphi^m}{m!}.
\end{equation}
On an all-time trajectory where $\mathcal M$ is bounded
below, $\int_0^\infty\mathcal C(t)\,dt<\infty$.
There are therefore times $t_j\to\infty$ with
$\mathcal C(t_j)\to0$. This says nothing by itself
about smooth compactness of the metrics.

Using only integrability, even pointwise decay of
$\mathcal C(t)$ requires additional temporal control:
a smooth nonnegative integrable function may have
arbitrarily narrow spikes of fixed height. Uniform
continuity excludes this. A fixed positive value
at infinitely many late times would then force
disjoint intervals of uniform length on which
the function is at least half as large, contrary
to integrability. This elementary argument
explains one use of uniform geometric estimates.

\paragraph{A conditional convergence package.}
Assume that an all-time solution has normalized potentials
precompact in $C^\infty$, with uniformly positive metrics,
and that $\mathcal M$ is bounded below on it. At the times
above, take a subsequence converging smoothly to
$\varphi_\infty$. Continuity of scalar curvature and
volume gives $S_{\varphi_\infty}=\overline S$.
This proves conditional subsequential convergence,
not the full target.

For a precise upgrade, put
$E(t)=\mathcal M(\varphi(t))-\mathcal M(\varphi_\infty)$.
Monotonicity and continuity along the subsequence give
$E(t)\downarrow0$. Suppose for all sufficiently large $t$
\begin{equation}\label{a401:gradient}
 \sqrt{\mathcal C(t)}\geq cE(t)^{1-\theta},
                 \qquad c>0,\quad0<\theta<1.
\end{equation}
Then, wherever $E>0$,
\[
 -\frac{d}{dt}E^\theta
 =\theta E^{\theta-1}\mathcal C
                      \geq c\theta\sqrt{\mathcal C}.
\]
Thus $\int_T^\infty\sqrt{\mathcal C(t)}\,dt<\infty$.
If $E$ reaches zero in finite time, the dissipation
identity makes the later metric stationary.

Uniform equivalence of volume forms, together
with subtraction of the fixed-volume mean,
bounds the normalized potential speed by
$\|\partial_t\varphi\|_{L^2(\omega)}
 \leq C\sqrt{\mathcal C(t)}$.
The potentials are consequently Cauchy in $L^2$.
Their assumed $C^\infty$ precompactness makes
every smooth cluster point equal to that
unique $L^2$ limit. Every sequence has such
a subsequence, so the whole flow converges
smoothly to the cscK metric.

Inequality \eqref{a401:gradient} is an additional
gradient inequality, including its validity
along the whole late-time trajectory. Neither
a global inequality, entry into its neighborhood,
nor control of holomorphic automorphism directions
is assumed to have been proved. Convergence only
after time-dependent automorphisms would also
be weaker than the fixed-metric conclusion here.
The value of this package is that every missing
analytic hypothesis can be identified separately.

\paragraph{The linearized equation has the correct sign.}
At a flat complex torus, scalar-curvature variation
is $DS(u)=-\Delta^2u$, using the complex Laplacian.
Hence the linearized flow is
\[
                         \partial_tu=-\Delta^2u.
\]
If $-\Delta e_k=\lambda_ke_k$, Fourier modes evolve as
$u_k(t)=e^{-\lambda_k^2t}u_k(0)$. Nonconstant modes
decay and the constant mode is the potential gauge.
This verifies linear stability and fourth-order
parabolicity, but not nonlinear global existence.

Fourth order prevents a direct second-order
maximum-principle argument. On a circle, consider
the scalar biharmonic heat equation and
$u_0(x)=(1-\cos x)^2\geq0$. At zero,
$u_0(0)=0$ and $u_0^{(4)}(0)=6$, giving
$\partial_tu(0,0)=-6$. Positivity is immediately
lost. This is a test of the leading operator,
not an identification of Kähler positivity
with positivity of a potential. A metric
positivity argument needs additional structure.

\paragraph{Small energy does not bound curvature.}
Work in the flat class of a complex torus. In
a coordinate ball choose a smooth real bump
$\psi$ whose linearized scalar curvature is
not identically zero. For small $\varepsilon$ set
\[
 \varphi_\varepsilon(x)
             =\varepsilon^{7/2}\psi(x/\varepsilon),
\]
extended by zero. Its Hessian is
$O(\varepsilon^{3/2})$, so
$\omega+i\partial\bar\partial\varphi_\varepsilon$
is positive and uniformly close to the flat
metric. Its fourth derivatives have size
$\varepsilon^{-1/2}$. With the indicated
choice of bump, scalar curvature, and hence
the curvature tensor, attain a size bounded
below by a positive multiple of this quantity.

The nonlinear errors are lower order: products
of third derivatives are $O(\varepsilon)$,
as is a second derivative multiplying a
fourth derivative. The volume of the support
in real dimension $d=2m$ is
$O(\varepsilon^d)$. The class average is zero,
and the metric-volume correction tends
uniformly to zero. Consequently
\[
 \mathcal C(\varphi_\varepsilon)
                =O(\varepsilon^{d-1})\longrightarrow0,
 \qquad \sup|\operatorname{Rm}|\longrightarrow\infty.
\]
These are smooth metrics in one fixed class
containing a cscK metric. They are not claimed
to lie on one Calabi trajectory or to cause
a finite-time singularity. They refute the
specific a priori implication from small
Calabi energy alone to bounded curvature.

The example also explains why smoothness of
each initial condition cannot be replaced
by a uniform estimate over all initial
conditions using just this energy. The
small support allows high derivatives to
escape an integral scalar-curvature bound.
Higher regularity, nonconcentration or an
independent compactness mechanism is needed.

\paragraph{Verification, first gap, and outcome.}
The derivative orders, bump exponents and
Fourier decay were checked independently.
The convergence deduction uses exactly the
displayed compactness and gradient hypotheses.
Finite-time continuation would need uniform
positivity and sufficient higher spatial
bounds to invoke fourth-order parabolic
regularity; the energy identity does not
provide them.

The first general gap is to exclude metric
degeneration or curvature concentration for
arbitrary initial metrics in the class.
Even after all-time existence, full
convergence needs a mechanism excluding
multiple cluster points and uncontrolled
drift. The stated package identifies one
such sufficient mechanism.
\textbf{UNRESOLVED.} Dissipation consequences,
conditional convergence and explicit
failure tests are proved, but neither
part of the unrestricted target is
established or claimed as new.

\subsection{Sources}
\begin{itemize}
\item[S1]Benjamin Sibley, Long-Time Existence for the Calabi Flow on Ruled Manifolds over Riemann Surfaces, Conjectures 1.1-1.2 (2022 preprint).
\url{https://geometry.ulb.ac.be/wp-content/uploads/2022/04/Calabi-Flow-Version-VII.pdf}.
\textit{Formulation source. Consulted 24 September 2026: Calabi-flow formulation and the constant-scalar-curvature convergence branch.}
\end{itemize}
\end{document}
